\documentclass[10pt]{article}
\usepackage[margin=1.55cm]{geometry}
\usepackage{graphicx}
\usepackage{amsmath}
\usepackage{amssymb}
\usepackage[hidelinks]{hyperref}
\setlength{\parskip}{2pt}
\setlength{\parindent}{0pt}
\graphicspath{{../analysis_output/accessibility_alpha_v2w4wd/}}

\title{\vspace{-1.4cm}Matched $V_2$ and $d4$-readout decay closes the scale escape\\
but erases the compositional circuit\\
\large 8 $\alpha$ $\times$ 8 $\lambda$ $\times$ 3 seeds = 192 cells}
\author{}
\date{\vspace{-1.2cm}}

\begin{document}
\maketitle
\vspace{-1.9cm}

\section*{Design}
This control is identical to the V2-only-decay experiment except that the same phase-2
coefficient $\lambda$ penalizes both $V_2$ and the $d4$ readout row $W[2]$.  The embedding,
$V_1$, and the two $d2$ readout rows remain decay-free; all parameters remain trainable and
the two $d2$ losses remain active.  The initial accessibility is set by
$V_1\mapsto\alpha V_1$ and measured as $R_0=Q_2/Q_1$.  The grid uses eight $\alpha$ values,
eight $\lambda$ values from 0 to $10^{-1}$, and three pretrained seeds.  All 192 cells
completed, retained $d4$ accuracy, and are represented in the trajectory files.

\section*{Findings}
\textbf{1. Readout decay does not sharpen the transition at
$R_c=1/\sqrt8$.}
Very weak decay gives a small increase in composition, but the effect reverses before a
phase develops.  At $R_0=0.417$, the mean $2{+}2$ fraction is
$0.101\;(\lambda=0)\to0.124\;(3{\cdot}10^{-4})$, then falls to
$0.015\;(10^{-3})$ and $0.0028\;(3{\cdot}10^{-3})$.  At $R_0=0.515$ the corresponding
sequence is $0.103\to0.128\to0.024\to0.0021$.  Pooling all supercritical cells gives
$f_{2+2}=0.112$ at $3{\cdot}10^{-4}$, $0.020$ at $10^{-3}$, and $0.0021$ at
$3{\cdot}10^{-3}$.  The phase diagrams therefore become uniformly blue as $\lambda$ grows;
there is no decay-sharpened boundary.

\textbf{2. The contrast with both benchmarks is large.}
At $R_0=0.417,\lambda=3{\cdot}10^{-3}$, matched readout decay gives
$f_{2+2}=0.0028$, compared with $0.185$ under V2-only decay.  At $\lambda=10^{-2}$ the
comparison is $0.0007$ versus $0.159$.  The frozen benchmark is more strongly
compositional still: around matched $R\simeq0.436$, it gives $f_{2+2}=0.84$--$0.89$ at
$3{\cdot}10^{-3}$ and $0.87$--$0.92$ at $10^{-2}$.  Readout decay thus does not complete
the V2-only experiment's partial transition; it removes it.

\begin{figure}[h!]
\centering
\includegraphics[width=0.49\textwidth]{phase_R0.pdf}\hfill
\includegraphics[width=0.49\textwidth]{phase_Rfinal.pdf}
\vspace{-4pt}
\caption{Mean compositional fraction versus $R_0$ (left) and $R_{\rm final}$ (right).
The red line marks $R_c$.  Weak decay preserves a faint accessibility dependence, but
$\lambda\geq10^{-3}$ suppresses composition on both sides of the theoretical boundary.}
\label{fig:phase}
\end{figure}

\textbf{3. $R(t)$ remains high while composition disappears.}
Decay raises $R_{\rm final}$ relative to the no-decay trajectory and compresses the sweep
toward a supercritical band.  For example, $R_0=0.323$ ends at $R_{\rm final}=0.311$ with
no decay but at $0.407$ for $\lambda=10^{-2}$; $R_0=0.515$ ends at $0.412$ and $0.443$,
respectively.  High-$R_0$ conditions can drift downward in absolute terms, while low-$R_0$
conditions drift upward.  The important observation is the decoupling: at
$\lambda=10^{-2}$ the representation is mostly above $R_c$, yet the compositional fraction
is only $4{\cdot}10^{-4}$ in the initially supercritical group.

\clearpage
\begin{figure}[h!]
\centering
\includegraphics[width=0.49\textwidth]{cuts_fixed_lambda.pdf}\hfill
\includegraphics[width=0.49\textwidth]{drift.pdf}
\vspace{-4pt}
\caption{Left: fixed-$\lambda$ cuts show the small weak-decay separation and its collapse
at larger $\lambda$.  Right: $R_{\rm final}$ versus $R_0$; decay funnels the representation
toward $R\approx0.4$--$0.45$ rather than restoring a compositional phase.}
\label{fig:drift}
\end{figure}

\textbf{4. The V2-only scaling escape is fixed, but balanced norms do not emerge.}
In the V2-only experiment, $V_2$ could vanish while the unpenalized readout stayed large:
at $R_0=0.417,\lambda=10^{-2}$, $\|V_2\|=4.7{\cdot}10^{-4}$ while the total readout norm
was $2.83$.  With matched decay, both relevant norms shrink: at the same $R_0$,
$(\|V_2\|,\|W_4\|)=(0.420,0.734)$ at $10^{-3}$,
$(0.116,0.372)$ at $10^{-2}$, and $(0.067,0.276)$ at $3{\cdot}10^{-2}$.  However, the
ratio $\|W_4\|/\|V_2\|$ grows from $1.75$ to $3.22$ to $4.11$ rather than approaching the
near-balanced frozen solution (roughly $1.8$ versus $1.9$ around the same $R$).  Matched
parameter decay therefore removes the literal rescaling loophole but does not select the
frozen compositional optimum in the full network.

\textbf{5. The first-block residual shortcut still blocks a pure compositional phase.}
For $\lambda\leq3{\cdot}10^{-2}$, $|\hat f_4|$ remains at least $0.998$ even where $V_2$
is tiny and $f_{2+2}$ is nearly zero.  Hundreds of block-2 neurons remain classified as
relevant, but their computation is not parent-compositional.  Because the unpenalized,
trainable first block feeds the readout through the residual stream, the most consistent
interpretation is that its learned $d4$ leak carries the target while decay dismantles the
two-parent block-2 route.  This is an inference from the present observables rather than a
direct route decomposition; the parallel-$d4$-adapter experiment records first- and
second-block contributions explicitly.

\clearpage
\begin{figure}[h!]
\centering
\includegraphics[width=0.98\textwidth]{phase_plane.pdf}
\vspace{-4pt}
\caption{Seed-0 trajectories in $(R(t),f_{2+2}(t))$.  Under appreciable matched decay,
trajectories move toward or remain in the supercritical-$R$ region while falling toward
zero composition.  Accessibility alone no longer predicts which circuit carries $d4$.}
\label{fig:plane}
\end{figure}

\textbf{Bottom line.}
Readout decay answers the scale-balance objection to V2-only decay, but the result is not a
sharper transition.  Instead, it reveals a stronger bypass: the full network keeps $d4$
solved, maintains or increases measured accessibility, and abandons the $2{+}2$ circuit.
The frozen benchmark's clean transition requires not only matched regularization but also
protection from alternative trainable routes.  Directly regularizing a new block-1 $d4$
adapter alongside $V_2$ and $W_4$, while freezing the pretrained $d2$ parents, is therefore
the appropriate next comparison.

\vspace{2pt}
\footnotesize
Artifacts: \texttt{analysis\_output/accessibility\_alpha\_v2w4wd/}
(merged results, trajectories, phase diagrams, cuts, drift, and phase-plane plot).  Code:
\texttt{analysis/accessibility/run\_alpha\_unfrozen\_experiment.py}
(\texttt{--wd-block2-readout}) and \texttt{scripts/launch\_v2w4wd\_workers.sh}.

\end{document}
